% IN5340/IN9340 Project II Presentation
% Build (from this folder): pdflatex -interaction=nonstopmode -halt-on-error presentation_p2.tex

\documentclass[aspectratio=169,10pt]{beamer}

\usetheme{Madrid}
\usecolortheme{default}

\usepackage{graphicx}
\usepackage{amsmath,amssymb}
\usepackage{booktabs}
\usepackage{mathtools}
\usepackage[T1]{fontenc}
\usepackage[utf8]{inputenc}
\usepackage{hyperref}

\graphicspath{{figs/}}

\newcommand{\fullfig}[2][]{%
\centering
\includegraphics[width=\linewidth,height=0.78\textheight,keepaspectratio,#1]{#2}%
}
\newcommand{\colfig}[2][]{%
\includegraphics[width=\linewidth,height=0.58\textheight,keepaspectratio,#1]{#2}%
}

\title[IN5340 Project II]{IN5340/IN9340 Project II: Estimation Theory}
\author{Theodor W\aa lberg}
\institute{University of Oslo}
\date{February 2026}

\begin{document}

\begin{frame}
\titlepage
\vspace{-1.0em}
\centering
\footnotesize
Code: \href{https://github.com/Twaal/Statistical_Signal_Processing_IN5340/tree/master/Project_2}{\nolinkurl{github.com/Twaal/Statistical_Signal_Processing_IN5340/tree/master/Project_2}}
\end{frame}

\begin{frame}{Agenda}
\begin{itemize}
\item Introduction and dataset
\item 1A-1D: Maximum likelihood time-delay estimation
\item 2A: Cram\'er-Rao lower bound
\item 3A-3B: Least squares delay estimation and comparison
\item Summary
\end{itemize}
\end{frame}

\section{Introduction}

\begin{frame}{Project scope and data}
\begin{itemize}
\item Goal: estimate time delay of strongest reflector in sonar data (32 channels).
\item Dataset: \texttt{sonardata2.mat}, complex matrix $\mathbf{data}\in\mathbb{C}^{8192\times 32}$.
\item Known parameters from assignment/code: $B=30\,$kHz, $f_c=100\,$kHz, $f_s=40\,$kHz, $T_p=8\,$ms.
\item Methods covered: ML (matched filter), CRLB, and LSE fine delay estimation.
\end{itemize}
\end{frame}

\begin{frame}{Raw channel example}
\fullfig{plot_01.png}
\end{frame}

\section{1. Maximum Likelihood Time Delay}

\subsection{Exercise 1A}

\begin{frame}{1A: ML estimator derivation}
Assume
\[
x(t)=s(t-T_0)+w(t),\qquad x[n]=s(n\Delta-\tau_0)+w[n],\; w[n]\sim\mathcal{N}(0,\sigma_w^2).
\]
For fixed $\tau_0$, maximizing likelihood is equivalent to minimizing
\[
J(\tau_0)=\sum_{n=0}^{N-1}\left(x[n]-s(n\Delta-\tau_0)\right)^2.
\]
Expanding $J(\tau_0)$ gives ML as cross-correlation maximization:
\[
\hat{\tau}_{\mathrm{ML}}=\arg\max_{\tau_0}\sum_{n=0}^{N-1}x[n]s(n\Delta-\tau_0).
\]
\vspace{0.5em}
\small
Interpretation: delay is chosen at matched-filter peak.
\end{frame}

\subsection{Exercise 1B}

\begin{frame}[fragile]{1B: Implementing chirp and matched filter}
\scriptsize
\begin{verbatim}
# chirp
B = 30000; fs = 40000; tp = 0.008
t = np.arange(0, tp/2, 1/fs)
a = B / tp
chirp = np.exp(1j * (2*np.pi) * a * t**2 / 2)

# matched filter for channel 16
x = data[:, 15]
matched_filter_output = correlate(x, chirp, mode='full')
\end{verbatim}
\normalsize
Positive-delay part is used.
\end{frame}

\begin{frame}{1B: Chirp signal (time and frequency)}
\fullfig{plot_02.png}
\end{frame}

\begin{frame}{1B: Matched filter output, channel 16}
\fullfig{plot_03.png}
\vspace{0.25em}
\small
Notebook print: correlation length $=8351$, original signal length $=8192$.
\end{frame}

\subsection{Exercise 1C}

\begin{frame}[fragile]{1C: Pulse compression on all channels}
\scriptsize
\begin{verbatim}
pc_list = []
for ch in range(num_channels):
    y_full = correlate(data[:, ch], chirp, mode='full')
    pc_list.append(y_full)
pulse_compressed_full = np.column_stack(pc_list)

start_idx = len(chirp) - 1
pulse_compressed = pulse_compressed_full[start_idx:, :]

before_mag_all = np.abs(data[:1600, :]).T
after_mag_all  = np.abs(pulse_compressed[:1600, :]).T
\end{verbatim}
\end{frame}

\begin{frame}{1C: Before vs after pulse compression}
\fullfig{plot_04.png}
\end{frame}

\begin{frame}{1C: Discussion}
\begin{itemize}
\item Before compression: energy is spread in time and appears more noise-like.
\item After matched filtering: peaks are sharper and stronger.
\item Detectability and time-delay resolution improve.
\item Output remains complex because both received data and chirp reference are complex.
\end{itemize}
\end{frame}

\subsection{Exercise 1D}

\begin{frame}[fragile]{1D: Peak detection per channel}
\scriptsize
\begin{verbatim}
peak_indices = np.argmax(after_mag_all, axis=1)

after_mag_all_upsampled_4 = resample(after_mag_all,
                                     after_mag_all.shape[1]*4, axis=1)
after_mag_all_upsampled_8 = resample(after_mag_all,
                                     after_mag_all.shape[1]*8, axis=1)
peak_indices_upsampled[4] = np.argmax(after_mag_all_upsampled_4, axis=1)
peak_indices_upsampled[8] = np.argmax(after_mag_all_upsampled_8, axis=1)
\end{verbatim}
\end{frame}

\begin{frame}{1D: Peak per channel, no upsampling}
\fullfig{plot_05.png}
\end{frame}

\begin{frame}{1D: Delay estimates with upsampling x4 and x8}
\fullfig{plot_06.png}
\end{frame}

\begin{frame}{1D: Interpretation}
\begin{itemize}
\item Delays are very similar across channels, with small channel-dependent differences.
\item Main peak location is around sample index $\approx 304$ in the 1600-sample window.
\item After upsampling, index values scale with factor, while physical delay stays consistent.
\end{itemize}
\end{frame}

\section{2. Cram\'er-Rao Lower Bound}

\subsection{Exercise 2A}

\begin{frame}[fragile]{2A: SNR and CRLB computation}
\scriptsize
\begin{verbatim}
peak_magnitudes = after_mag_all[np.arange(after_mag_all.shape[0]), peak_indices]
peak_intensities = peak_magnitudes**2
noise_intensities = np.mean(np.delete(after_mag_all**2, peak_indices, axis=1), axis=1)
snr = peak_intensities / noise_intensities

b2_rms = (2*np.pi)**2 * ((fc**2) + (B**2)/12)
crlb = 1 / (2 * B * tp * snr * b2_rms)
\end{verbatim}
\small
Notebook print: RMS bandwidth approximation $\beta_{\mathrm{rms}}^2 \approx 3.977\times 10^{11}$.
\end{frame}

\begin{frame}{2A: CRLB across channels}
\fullfig{plot_07.png}
\end{frame}

\begin{frame}{2A: Interpretation}
\begin{itemize}
\item CRLB values are not identical for all channels.
\item They remain in the same order of magnitude.
\item Differences are expected due to finite samples and channel-dependent noise/reflector levels.
\end{itemize}
\end{frame}

\section{3. Least Squares Estimation}

\subsection{Exercise 3A}

\begin{frame}{3A: LSE procedure}
For pulse-compressed intensity $I[n]=|x_{pc}[n]|^2$ in the first 1600 samples:
\begin{enumerate}
\item Upsample with factor 4 or 8.
\item Locate dominant peak index.
\item Keep samples around peak above half-power threshold.
\item Form $y=\ln(I)$.
\item Build observation matrix $H=[\mathbf{1},\;t,\;t^2]$.
\item Solve LS for $\ln(I(t))\approx a+bt+ct^2$ and get fine delay from vertex.
\end{enumerate}
\end{frame}

\begin{frame}[fragile]{3A: LS model and fine-delay formula}
\scriptsize
\begin{verbatim}
y = np.log(upsampled_intensities[ch, selected_indices])
t_selected = selected_indices * dt
H = np.column_stack([np.ones(len(t_selected)), t_selected, t_selected**2])
a_hat, b_hat, c_hat = np.linalg.lstsq(H, y, rcond=None)[0]
t_peak = -b_hat / (2 * c_hat)
\end{verbatim}
\normalsize
Model:
\[
\ln I(t) = a + bt + ct^2, \qquad \hat\tau_{\mathrm{LSE}} = -\frac{\hat b}{2\hat c}.
\]
\end{frame}

\begin{frame}{3A: Original vs upsampled intensity around peak, channel 16}
\fullfig{plot_08.png}
\end{frame}

\begin{frame}{3A: Log-intensity and quadratic fit}
\fullfig{plot_09.png}
\end{frame}

\subsection{Exercise 3B}

\begin{frame}{3B: Delay estimates across channels, ML vs LSE}
\fullfig{plot_10.png}
\end{frame}

\begin{frame}{3B: Difference between LSE and upsampled ML}
\fullfig{plot_11.png}
\end{frame}

\begin{frame}{3B: Discussion}
\begin{itemize}
\item Delays are close across channels, but not exactly identical.
\item LSE gives smoother sub-sample continuous estimates.
\item ML argmax is constrained to the sample grid after optional upsampling.
\item LSE and ML track the same reflector, with small residual differences due to noise and local pulse-shape variation.
\end{itemize}
\end{frame}

\section*{Summary}

\begin{frame}{Summary}
\begin{itemize}
\item Exercise 1A-1D: Derived and implemented ML delay estimation via matched filtering; validated peak consistency across channels and upsampling factors.
\item Exercise 2A: Computed CRLB per channel from estimated SNR; bounds are similar but channel-dependent.
\item Exercise 3A-3B: Implemented LSE fine delay from log-intensity quadratic fit; obtained sub-sample delay refinement and compared against ML.
\item All figures in this deck are exported directly from executed \texttt{project\_2.ipynb} outputs.
\end{itemize}
\end{frame}

\end{document}
